\begin{textsample}{POS Dim 3 – ac – Score 53.00 – ac/ac\_inf\_20.txt}  \label{ex:f3_pos_023}
20. \textbf{CRIANÇAS} BEM \textbf{PEQUENAS} ( 1 ANO E 7 \textbf{MESES} A 3 ANOS E 11 \textbf{MESES} ) 20.1.
Campo de experiências “ \textbf{Escuta}, fala, pensamento e imaginação ” Desde o \textbf{nascimento}, as \textbf{crianças} participam de situações comunicativas cotidianas com as pessoas com as quais interagem.
As primeiras formas de interação do \textbf{bebê} são os movimentos do seu corpo, o olhar, a postura corporal, o sorriso, o choro e outros recursos vocais, que ganham sentido com a interpretação do outro.
Progressivamente, as \textbf{crianças} vão ampliando e enriquecendo seu vocabulário e demais recursos de expressão e de compreensão, apropriando -se da língua materna – que se torna, pouco a pouco, seu veículo privilegiado de interação.
Na Educação \textbf{Infantil}, é importante promover experiências nas quais as \textbf{crianças} possam falar e \textbf{ouvir}, potencializando sua participação na cultura oral, pois é na \textbf{escuta} de histórias, na participação em conversas, nas descrições, nas narrativas elaboradas individualmente ou em grupo e nas implicações com as múltiplas linguagens que a \textbf{criança} se constitui ativamente como sujeito singular e pertencente a um grupo social.
Desde cedo, a \textbf{criança} manifesta \textbf{curiosidade} com relação à cultura escrita : ao \textbf{ouvir} e acompanhar a leitura de textos, ao observar os muitos textos que circulam no contexto familiar, comunitário e escolar, ela vai construindo sua concepção de língua escrita, reconhecendo diferentes usos sociais da escrita, dos gêneros, suportes e portadores.
Na Educação \textbf{Infantil}, a \textbf{imersão} na cultura escrita deve partir do que as \textbf{crianças} conhecem e das \textbf{curiosidades} que deixam transparecer.
As experiências com a literatura \textbf{infantil}, propostas pelo educador, mediador entre os textos e as \textbf{crianças}, contribuem para o desenvolvimento do gosto pela leitura, do \textbf{estímulo} à imaginação e da ampliação do conhecimento de mundo.
Além disso, o contato com histórias, contos, fábulas, poemas, cordéis etc. propicia a familiaridade com \textbf{livros}, com diferentes gêneros literários, a diferenciação entre ilustrações e escrita, a aprendizagem da direção da escrita e as formas corretas de manipulação de \textbf{livros}.
Nesse convívio com textos escritos, as \textbf{crianças} vão construindo hipóteses sobre a escrita que se revelam, inicialmente, em rabiscos e garatujas e, à medida que vão conhecendo letras, em escritas espontâneas, não convencionais, mas já \textbf{indicativas} da compreensão da escrita como sistema de representação da língua. 20.2.
Direitos de aprendizagem e desenvolvimento - Conviver com \textbf{crianças} e \textbf{adultos} em situações comunicativas do cotidiano, desenvolvendo progressivamente confiança em participar de experiências que envolvam o pensamento, a imaginação, sentimentos, narrativas e diálogos, quer seja em momentos individuais ou coletivos. - \textbf{Brincar} e se deliciar com histórias e poesias, com jogos, parlendas, rimas, \textbf{brinquedos} cantados, textos de imagens e escritos, ampliando seu repertório das manifestações culturais, favorecendo sua progressiva aprendizagem dos vários gêneros e formas de expressão : gestual, oral, escrita, \textbf{plástica}, \textbf{dramática} e musical. - Explorar \textbf{gestos}, expressões, rimas, imagens, canções, enredos de histórias, textos que compõem o patrimônio cultural da humanidade ( histórias, lendas, mitos, fábulas, poemas, parlendas, trava-línguas, piadas, adivinhas, músicas, entre outros ), atuando como autores no processo, redescobrindo e transformando à sua maneira e de acordo com suas possibilidades, considerando o momento de seu desenvolvimento. - Participar de \textbf{rodas} de conversa, relatos de experiência, leitura e contação de histórias, saraus, dramatizações, apresentações, construção de narrativas, representação de papéis no faz de conta, exploração de materiais impressos e de variedade linguística, comunicando -se com os outros, trocando informações e agindo no mundo, criando assim a possibilidade de compreensão, de organização de ideias e consciência de si mesmo. - Expressar sentimentos, ideias, \textbf{desejos}, necessidades, dúvidas, informações e descobertas, possibilitando o compartilhamento de saberes e conhecimentos entre todos os envolvidos, sendo sua autoria garantida pelo \textbf{gesto}, pela fala, pelo desenho, pela escrita ( convencional ou não convencional ) ou por outra forma de \textbf{registro}. - Conhecer -se como sujeitos de sua cultura, reconhecendo e dando sentido às intenções comunicativas e de \textbf{escuta}, aos \textbf{gestos}, às emissões sonoras, às narrativas de diferentes gêneros linguísticos, aos diferentes suportes e gêneros textuais, valorizando suas origens. 20.2.1.
Objetivo de aprendizagem e desenvolvimento : ( EI02EF01 ) - Dialogar com \textbf{crianças} e \textbf{adultos}, expressando seus \textbf{desejos}, necessidades, sentimentos e opiniões.
Aprendizagens específicas - Expressar e comunicar ideias, \textbf{desejos}, necessidades, sentimentos e vivências por meio de diferentes linguagens ( dança, desenho, \textbf{mímica}, música, linguagem oral e escrita e outras ). - Valorizar a conversa como forma de conhecer a si mesma, as coisas e aos outros. - Expressar suas emoções e sentimentos de modo que seus \textbf{hábitos}, ritmos e preferências individuais sejam respeitadas no grupo em que convive. - Desenvolver a \textbf{escuta} \textbf{atenta} nas interações comunicativas com seus \textbf{pares} e \textbf{adultos}. - Ampliar, progressivamente, o vocabulário acrescentando palavras às construções de frases. - \textbf{Interessar} -se por fazer uso mais complexo da linguagem, na tentativa de se fazer entender. - Desenvolver o pensamento e o raciocínio para compreender e expressar livremente os \textbf{sons} da língua e a sonoridade das palavras. - Desenvolver, progressivamente, habilidade de iniciar diálogos estruturados. \textbf{Sugestões} de experiências - \textbf{Conversar} com os \textbf{colegas} e com o professor sobre situações vividas no cotidiano. - Relatar suas próprias experiências. - \textbf{Escutar} as experiências vividas pelos \textbf{colegas}. - Falar sobre seus estados emocionais : satisfação e insatisfação, prazer e desagrado, surpresa, temor ou preocupação. - Cumprimentar, despedir -se, agradecer, desculpar -se, aceitar desculpas. - Dançar, fazer \textbf{mímica}, desenhar, cantar e \textbf{dramatizar}. - \textbf{Ouvir} \textbf{contar} e recontar histórias, parlendas, quadrinhas, trava-línguas, adivinhas e poesias. - Nomear e descrever objetos, pessoas, fotografias e gravuras. - \textbf{Emitir} \textbf{sons} articulados a \textbf{gestos} e ser interpretado pelo grupo de \textbf{crianças} e pelo \textbf{adulto}. 20.2.2.
Objetivo de aprendizagem e desenvolvimento : ( EI02EF02 ) - Identificar e criar diferentes \textbf{sons} e reconhecer rimas e aliterações em cantigas de \textbf{roda} e textos poéticos.
Aprendizagens específicas - Conhecer, reproduzir e produzir oralmente jogos verbais como trava-línguas, parlendas, adivinhas, quadrinhas, poemas e canções. - Ampliar seu vocabulário por meio de músicas, cantigas de \textbf{roda}, narrativas e \textbf{brincadeiras}, desenvolvendo sua capacidade de comunicação. - Desenvolver o interesse por explorar \textbf{sons}, seus efeitos e intensidades, nos jogos rimados. - Desenvolver a imaginação e a criatividade, \textbf{brincando} com a linguagem em canções conhecidas, danças, recitações de poemas, parlendas, cantigas de \textbf{roda}, dentre outras. - Construir diferentes entonações, ritmos e \textbf{sons}, enquanto canta ou declama. - Desenvolver o \textbf{hábito} da \textbf{escuta} ao recontar e usar em suas \textbf{brincadeiras}, rimas e textos poéticos. - Familiarizar -se, gradativamente, com aspectos da língua por meio da musicalidade e forma gráfica dos textos poéticos e rimados. - Desenvolver o gosto e o prazer em realizar suas produções e declamações. \textbf{Sugestões} de experiências - Recitar poesias e parlendas criando diferentes entonações e ritmos. - Declamar textos poéticos conhecidos. - Participar de jogos orais envolvendo \textbf{sons}, ritmos e rimas. - Criar \textbf{sons} enquanto canta uma música. - \textbf{Imitar} a leitura de textos poéticos, trechos de contos, listas e outras possibilidades. - \textbf{Brincar} com palavras que rimam. - Recitar poemas e/ou cantar músicas com rimas. - \textbf{Brincar} de trava-línguas, parlendas, adivinhas, quadrinhas, poemas e canções. - Jogar e \textbf{brincar} com a sonoridade das palavras. 20.2.3.
Objetivo de aprendizagem e desenvolvimento : ( EI02EF03 ) - \textbf{Demonstrar} interesse e atenção ao \textbf{ouvir} a leitura de histórias e outros textos, diferenciando escrita de ilustrações, e acompanhando, com orientação do adulto-leitor, a direção da leitura ( de cima para baixo, da esquerda para a direita ).
Aprendizagens específicas - Ampliar, progressivamente, a capacidade de manter -se \textbf{atento} durante a leitura do professor. - Perceber que as ilustrações de um \textbf{livro} podem representar o que está escrito na narrativa ou agregar novas informações ao texto. - Desenvolver a compreensão de textos, a partir da \textbf{escuta} e interpretação dos \textbf{sons}, \textbf{gestos} e expressões que fazem parte das narrativas. - Apropriar -se, gradativamente, da linguagem escrita ( nos diferentes tipos de textos ), sua função, conteúdo e formato. - Desenvolver, progressivamente, comportamento leitor, observando \textbf{gestos} e ações que acompanham a leitura de um \textbf{livro}. - Apropriar -se, progressivamente, dos procedimentos leitores pela imitação : \textbf{gestos} de folhear, apontar para as palavras e imagens, buscar títulos no índice, ler da esquerda para direita e de cima para baixo, dentre outros. - Produzir ilustrações estabelecendo relação entre os fatos da narrativa. - Desenvolver a compreensão da narrativa, por meio da produção de ilustrações sequenciada dos fatos. - Construir texto oral a partir das reflexões sobre as ilustrações de textos e \textbf{livros}. \textbf{Sugestões} de experiências - Ler textos memorizados com a \textbf{ajuda} do ( a ) professor ( a ). - \textbf{Brincar} de ler em diferentes situações. - \textbf{Brincar} de ler \textbf{livros}, revistas e outros materiais impressos a partir da capa e virar as páginas sucessivamente. - \textbf{Brincar} de ler acompanhando o texto com o \textbf{dedo}. - \textbf{Conversar} sobre as ilustrações dos \textbf{livros}, revistas, gibis e outros impressos. - \textbf{Conversar} sobre as histórias e outros textos lidos. - Nomear e descrever objetos, pessoas, fotografias e gravuras. - Apreciar a leitura de histórias feita pelo professor. - \textbf{Escutar} leituras diversas. - Criar ilustrações variadas. - \textbf{Contar} histórias a partir de objetos \textbf{lúdicos} e da imaginação. 20.2.4.
Objetivo de aprendizagem e desenvolvimento : ( EI02EF04 ) - Formular e responder perguntas sobre fatos da história narrada, identificando cenários, \textbf{personagens} e principais acontecimentos.
Aprendizagens específicas - Desenvolver, gradativamente, a compreensão sobre a estrutura da narrativa a partir de diversas reflexões. - Desenvolver, gradativamente, sua capacidade argumentativa, opinando sobre características de \textbf{personagens}, cenários e os principais acontecimentos de histórias conhecidas, lidas e narradas. - Desenvolver a capacidade de recuperação e ordenação da narrativa, por meio de reconto oral, dramatização e ilustrações. \textbf{Sugestões} de experiências - Descrever, oralmente, características de \textbf{personagens} e cenários. - \textbf{Conversar} com seus \textbf{pares} sobre as histórias conhecidas. - Recontar histórias \textbf{apoiadas} por ilustrações. - Representar histórias conhecidas e seus \textbf{personagens} por meio de dramatizações. - \textbf{Ouvir} e apreciar histórias e outros textos literários. 20.2.5.
Objetivo de aprendizagem e desenvolvimento : ( EI02EF05 ) - Relatar experiências e fatos acontecidos, histórias \textbf{ouvidas}, filmes ou peças teatrais assistidos etc.
Aprendizagens específicas - Desenvolver, gradativamente, a capacidade argumentativa e crítica, elaborando perguntas e respostas de acordo com as diversas situações comunicativas nas quais participa. - Expressar -se livremente, por meio da oralidade, ideias e sentimentos, respondendo e formulando perguntas. - Compartilhar oralmente suas experiências vividas e \textbf{ouvidas}, descrevendo lugares, pessoas e objetos. - Desenvolver a capacidade de construir narrativas orais. - Ampliar o vocabulário, participando e expressando -se em conversas, narrações e \textbf{brincadeiras}. - Compreender, gradativamente, o conteúdo e o propósito de diferentes mensagens em diversos contextos. - Desenvolver postura de respeito e \textbf{escuta} à fala do outro. \textbf{Sugestões} de experiências - Expressar -se verbalmente em conversas, narrações e \textbf{brincadeiras}. - Relatar experiências pessoais. - \textbf{Contar} e \textbf{ouvir} casos, relatos. - Pedir e atender pedidos, dar e \textbf{ouvir} recados, avisos, orientações e instruções. - Dar e \textbf{ouvir} notícias e/ou informações. - \textbf{Ouvir}, \textbf{contar} e recontar histórias, trava-línguas, parlendas, quadrinhas. - \textbf{Brincar} de faz de conta. - Participar oralmente da construção de \textbf{pequenos} textos coletivos que envolvam situações reais e de faz de conta. 20.2.6.
Objetivo de aprendizagem e desenvolvimento : ( EI02EF06 ) - Criar e \textbf{contar} histórias oralmente, com base em imagens ou temas sugeridos.
Aprendizagens específicas - Recontar histórias e contos conhecidos, com \textbf{ajuda} do professor, dos \textbf{colegas} e dos familiares. - Manifestar as experiências vividas e \textbf{ouvidas}, por meio de dramatizações ou outras representações. - Desenvolver a \textbf{autoconfiança} e segurança na capacidade comunicativa ao criar ou \textbf{contar} suas histórias. - Apropriar -se de recursos expressivos da narrativa de uma história, fazendo uso nas \textbf{brincadeiras}, representações e conversações com seus \textbf{pares} e \textbf{adultos}. - Desenvolver, progressivamente, a compreensão sobre o uso da linguagem e sua importância nas relações sociais. - Desenvolver a capacidade de estabelecer relações entre diferentes histórias conhecidas. \textbf{Sugestões} de experiências - Recontar histórias ao \textbf{brincar} de faz de conta. - Recontar histórias conhecidas com e sem apoio de \textbf{livros}, fantoches, desenhos, entre outros. - Ditar histórias criadas ou memorizadas. - Criar histórias com ou sem apoio de imagens, \textbf{fotos} e outros materiais. - \textbf{Brincar} de faz de conta construindo histórias. 20.2.7.
Objetivo de aprendizagem e desenvolvimento : ( EI02EF07 ) - \textbf{Manusear} diferentes portadores textuais, \textbf{demonstrando} reconhecer seus usos sociais.
Aprendizagens específicas - Familiarizar -se com o mundo letrado, realizando leitura não-convencional em diferentes portadores textuais. - Compreender, gradativamente, a função dos diferentes portadores textuais, fazendo uso no seu cotidiano. - Produzir escritas espontâneas em \textbf{brincadeiras} de faz de conta, tendo como base os diferentes portadores textuais. \textbf{Sugestões} de experiências - \textbf{Brincar} de escrever cartas aos seus \textbf{colegas} ou familiares, espontaneamente. - Folhear \textbf{livros}, \textbf{contando} suas histórias para os \textbf{colegas}. - \textbf{Brincar} de correio, escritório, supermercado, banco, livraria. - Escolher e selecionar \textbf{livros} para ler. - Explorar \textbf{livros} confeccionados com materiais diversos ( papel, \textbf{plástico}, \textbf{tecido}, borracha ). - Presenciar situações significativas de leitura e escrita. - Ler e escrever livremente o próprio nome e o nome dos \textbf{colegas}. - Visitar \textbf{bibliotecas} e casas de leituras. 20.2.8.
Objetivo de aprendizagem e desenvolvimento : ( EI02EF08 ) - \textbf{Manipular} textos e participar de situações de \textbf{escuta} para ampliar seu contato com diferentes gêneros textuais ( parlendas, histórias de aventura, tirinhas, cartazes de sala, cardápios, notícias etc. ).
Aprendizagens específicas - Familiarizar -se com os diferentes gêneros textuais em diferentes suportes. - Desenvolver atitude de \textbf{escuta}, em situações de conversação e de exploração dos diferentes gêneros textuais. - Desenvolver, gradativamente, a compreensão sobre as \textbf{marcas} e características dos diferentes gêneros textuais. - Desenvolver, gradativamente, a compreensão do uso social dos diferentes gêneros textuais. - Ampliar o gosto e o prazer pela leitura de diferentes gêneros textuais em diferentes suportes. \textbf{Sugestões} de experiências - Recitar parlendas. - Buscar informações nos jornais, revistas e outros. - Participar de atividades \textbf{culinárias}. - Usar \textbf{livros} de receitas, para preparações \textbf{culinárias}. - Explorar portadores de diferentes gêneros textuais. - \textbf{Manusear} suportes e gêneros textuais da nossa localidade. - Participar das situações de leitura de textos de diferentes gêneros ( contos, poemas, parlendas, trava-línguas, etc. ) feitas pelo professor. - Explorar, nas \textbf{rodas} de leitura, diferentes materiais escritos, tais como : \textbf{livros}, revistas, histórias em quadrinhos, dentre outros. - Participar de situações de escrita de avisos, recados, bilhetes e outros. - Ditar lista que fizerem sentido : frutas, \textbf{brinquedos}, \textbf{brincadeiras}, músicas e histórias que mais gostam. 20.2.9.
Objetivo de aprendizagem e desenvolvimento : ( EI02EF09 ) - \textbf{Manusear} diferentes instrumentos e suportes de escrita para desenhar, traçar letras e outros sinais gráficos.
Aprendizagens específicas - Desenvolver, gradativamente, a compreensão sobre as relações entre a fala e a escrita, a partir da observação do outro nos diversos contextos. - Compreender que desenhos, letras e outros sinais gráficos são formas de comunicação. - Compreender, gradativamente, o que a escrita representa e como é representada. - Desenvolver comportamento de escritor ao fazer de conta que escreve, \textbf{imitando} o \textbf{adulto}. - Desafiar -se a produzir, autonomamente, comunicações escritas não-convencionais ( garatujas ), \textbf{manuseando} diferentes instrumentos e suportes de escrita. \textbf{Sugestões} de experiências - \textbf{Manusear} e explorar diferentes suportes de escritas. - \textbf{Imitar} comportamento de escritor nas situações de faz de conta. - Desenhar livremente. - \textbf{Brincar} com a sonoridade das palavras. - \textbf{Brincar} de traçar letras e palavras. - Fazer uso de letras, ainda que de forma não convencional, em seus \textbf{registros} de comunicação. - Produzir comunicações escritas, tendo o professor como escriba. - Produzir diferentes \textbf{marcas} gráficas utilizando instrumentos e suportes variados. - Escrever o próprio nome – do jeito que souber. 20.3.
Avaliação - \textbf{Acompanhamento} do processo de aprendizagem e desenvolvimento das \textbf{crianças} bem \textbf{pequenas} Observação \textbf{criteriosa}, \textbf{registro} e análise quanto : - À capacidade de estabelecer interações comunicativas com seus \textbf{pares}, entendendo e se fazendo entender. - À capacidade de explorar \textbf{sons}, ritmos, entonações e intensidade nas rimas e aliterações. - À capacidade de estabelecer relação entre imagem e texto escrito. - Ao interesse em \textbf{ouvir} leitura de histórias e outros textos. - Aos procedimentos utilizados para \textbf{manusear} e \textbf{manipular} textos escritos de diferentes gêneros. - À capacidade argumentativa sobre elementos presentes em textos e fatos de histórias narradas. - Às iniciativas em relatar suas diversas experiências cotidianas. - À capacidade de criar histórias orais, estabelecendo relação entre imagem e temas propostos. - À compreensão dos usos e funções sociais da linguagem escrita. - À compreensão que desenhos, letras e outros sinais gráficos são formas de comunicação. - Às estratégias utilizadas para \textbf{manusear} diferentes instrumentos e suportes de escrita.

% matched lemmas: acompanhamento, adulto, ajuda, apoiar, atento, autoconfiança, bebê, biblioteca, brincadeira, brincar, brinquedo, colega, contar, conversar, criança, criterioso, culinário, curiosidade, dedo, demonstrar, desejo, dramatizar, dramático, emitir, escuta, escutar, estímulo, foto, gesto, hábito, imersão, imitar, indicativo, infantil, interessar, livro, lúdico, manipular, manusear, marca, mês, mímico, nascimento, ouvir, par, pequeno, personagem, plástico, registro, roda, som, sugestão, tecido
\end{textsample}
